\documentclass[a4paper,12pt]{report}
\usepackage[french]{babel}
\usepackage{enumitem}
\usepackage[utf8]{inputenc}% Pour utiliser les lettres accentuées
\usepackage{lscape}
\setlength{\voffset}{-0.5in}
\setlength{\footskip}{80pt}
\usepackage[T1]{fontenc}
\usepackage{color}
\usepackage{url}
\usepackage{natbib}
\usepackage{url}
\def\UrlBreaks{\do\/\do-}
\usepackage{breakurl}
\usepackage[breaklinks]{hyperref}
\bibliographystyle{unsrt}
\usepackage{graphicx}% utilisé les photos
\usepackage{xspace} %utilisé les espaces
%\usepackage{multirow}
%\usepackage{fancyhdr}
\usepackage[  bottom=2cm,left=2cm, right=2cm,]{geometry}
%\usepackage{soul}
\usepackage{array}
\newcommand{\cRM}[1]{\Roman {chapter}} % Capital
\newcommand{\cRm}[1]{\textsc{\romannumeral #1}} % Petit majuscule
\newcommand{\crm}[1]{\romannumeral #1}
% Siècle %
\newcommand{\siecle}[1]{\cRM{#1}\textsuperscript{e}~siècle}
\usepackage[nottoc]{tocbibind} 


\usepackage{fancyhdr}
\usepackage{amsthm}
\newtheorem{MaDefinition}{Définition}[section]
\setcounter{secnumdepth}{4}
\setcounter{tocdepth}{4} 
\usepackage{a4wide}
\usepackage{algorithm,algorithmic}
\usepackage{eurosym}
\usepackage{wrapfig}
\usepackage{graphicx}
\usepackage{booktabs}
\usepackage[table,xcdraw]{xcolor}
\usepackage{enumitem}
\usepackage{datatool}
\usepackage[normalem]{ulem}
\useunder{\uline}{\ul}{}

% Hauts et pieds de page
%\usepackage{fancyhdr}
%\pagestyle{fancy}
%\renewcommand{\headrulewidth}{1pt}
%\renewcommand{\footrulewidth}{1pt}



%\usepackage{fancyhdr} %pour les entêtes
%\pagestyle{fancy}
%*******************************************************************************************
\usepackage{array,multirow,makecell}
\setcellgapes{1pt}
\makegapedcells
\newcolumntype{R}[1]{>{\raggedleft\arraybackslash }b{#1}}
\newcolumntype{L}[1]{>{\raggedright\arraybackslash }b{#1}}
\newcolumntype{C}[1]{>{\centering\arraybackslash }b{#1}}



\pagestyle{fancy}
%definie l'entete se chaque page
\renewcommand\headrulewidth{1pt}
\fancyhf{}
\fancyhead[R]{\slshape \textbf \leftmark}
\fancyfoot[C]{\thepage}

%definition de l'entete du chaq chapitre
\makeatletter
\def\thickhrulefill{\leavevmode \leaders \hrule height 1ex \hfill \kern \z@}
\def\@makechapterhead#1{%
  %\vspace*{50\p@}%
  \vspace*{10\p@}%
  {\parindent \z@ \centering \reset@font

        \thickhrulefill\quad
        %\scshape 
        \@chapapp{} \thechapter
        \quad \thickhrulefill
        \par\nobreak
        \vspace*{10\p@}%
        \interlinepenalty\@M
        \hrule
        \vspace*{10\p@}%
        \Huge \bfseries #1\par\nobreak
        \par
        \vspace*{10\p@}%
        \hrule
    %\vskip 40\p@
    \vskip 100\p@
  }}

\renewcommand{\baselinestretch}{1.2}
%-----------------Début de la Page de garde---------------------------------
\renewcommand{\maketitle}{
\thispagestyle{empty}
\begin{titlepage}
\begin{center}
{\bf  République Algérienne Démocratique et Populaire
%\vspace{0.1cm}
\\ Ministère de l'Enseignement Supérieur et de la
Recherche Scientifique} \vspace{0.25cm}\\
 {\bf {\large Université A/Mira de Béjaïa}}\\
{\bf Faculté des Sciences Exactes} \\ {\bf Département d'Informatique}\\

%-------------------Pour insérer un logo
\begin{figure}[h!]
  \begin{center}
  % \usepackage{graphicx}
\includegraphics[width=4cm,height=2cm]{images/logo.jpg}\\
   \end{center}
\end{figure}
%-------------------------------------
\begin{center}
\vspace*{-1cm}
  \bf{MÉMOIRE DE MASTER } \\
 
  \bf{Option : Génie Logiciel } \\
\end{center}
\vspace{0.5cm}
{\LARGE {\bf Thème}} \\
\noindent\rule{\textwidth}{0.7mm}
\LARGE{\textbf{ Conception et réalisation d’une plateforme de gestion du marché du BTP en Algérie basée sur l’architecture REST
 } }
\noindent\rule{\textwidth}{0.7mm}
\end{center}
\begin{center}
\textbf{Réalisé par:}\\
\vspace{0.2cm}
%\textbf{ \hspace{0.1cm}$\textbf{M}$ \textsl{IDRISSOU}  Hamza}\\

\textbf{ \hspace{0.1cm}$\textbf{M.}$ \textsl{CHETTOUT}  Abdeslam}\\
\textbf{ \hspace{0.1cm}$\textbf{M.}$ \textsl{HAMMOUCHE}  Mohand Mouloud}\\

\end{center}

\vspace{0.1cm}
\textbf{Devant le jury composé de :}  \\
\begin{tabbing}
 \textbf{Examinateur 1}  \hspace{1.2cm}\= \textbf{M. -- --} \hspace{2.4cm}\=\textbf{} \\
 
 \textbf{Examinateur 2}    \> \textbf{M. --  --} \> \\ %\textbf{M.C.A}  \\
 \textbf{Encadrant}     \> \textbf{M. AKILAL Karim}  \> %\textbf{M.C.A}         \\
 \end{tabbing}
\newline
\newline

\newline
\begin{center}
\textbf{2021/2022}
\end{center}
\end{titlepage} 

}
%-----------------Fin de la Page de garde-----------------------------------------

%-----------------Début du document----------------------------------------------
\begin{document}
%-----------------Page de garde -----------------------------------------------
\maketitle
%--------------------------Abstract
\begin{abstract}
--Le Résumé--\\


\textbf{Mots clés:} \textit{----,----}
\vspace{0.2cm}
\begin{center}
\textbf{Abstract}
\end{center}
------\\

\textbf{Keywords: }\textit{MDM, MDS, SQL Server, ETL, DQS}


\end{abstract}
%--------------------------les remerciements
\chapter*{Remerciements}
\thispagestyle{empty}
\begin{it}
\begin{center}
  Nos remerciement vont en premier lieu  à nos parents qui ont su nous donner courage, volonté et patience afin d'achever ce projet ainsi tous les projets de la vie.\newline
  \newline
  Nous tenons à remercier également  \textbf{ Mr.AKILAL Karim} notre promoteur, pour son aide, ses conseils, ses encouragements et ses orientations qui nous ont guidés jusqu’à l’aboutissement de ce travail. \newline
  \newline
  
  Merci aux membres du jury d’avoir accepté de juger notre travail et de nous faire part de leurs remarques pertinentes.
\newline

Nous exprimons notre plus grande gratitude à tous ceux qui ont contribué de près ou de loin à la réalisation de ce travail.

  
\end{center}
\end{it}
\chapter*{Dédicaces}
\thispagestyle{empty}
\begin{it}

\begin{flushleft}
À-------- \newline
\newline

À--------\newline
\newline


À--------\newline
\newline

À--------\newline
\newline


À---------\newline
\newline

----------\newline
\begin{flushright}
\textbf{\textit{CHETTOUT Abdeslam}}
\end{flushright}
      

\end{flushleft}
\end{it}

%---------------------------Les dédicaces
%\include{Remerciementss}
\chapter*{Dédicaces}
\thispagestyle{empty}
\begin{it}

\begin{flushleft}
Je dédie ce modeste travail :\newline
\newline
À---------\newline
\newline

À----------\newline
\newline
À----------\newline
\newline

À----------\newline
\newline

----------


\begin{flushright}
\textbf{\textit{HAMMOUCHE Mohand Mouloud}}
\end{flushright}
\end{flushleft}
\end{it}

\pagenumbering{roman}
\setcounter{page}{1}
%------------------------  Pour l'entête
%\lhead{}
%--------------------------Tables des matières
\tableofcontents
\addtocontents{toc}{\protect\pagenumbering{roman}}  % pour la table des matières
\addcontentsline{toc}{chapter}{Table des matières}
%--------------------------liste des figures

\renewcommand\listfigurename{Liste des figures}
\listoffigures 
%--------------------------liste des tableaux
\listoftables
%--------------------------liste des algorithmes
%\newpage\addcontentsline{toc}{chapter}{Liste des algorithmes}
%\include{ListeDesAlgorithmes}
%--------------------------liste des abreviations
\newpage\addcontentsline{toc}{chapter}{Liste des abréviations}
\chapter*{Liste des abréviations}





\begin{description}
\parindent=2pt
 %\item[BDD :] Base de données.
 % \item[SDI :] \hspace{0.6cm}Système d'infomation décisionnelle.
  %\item[SGBD:] \hspace{0.1cm}Système de Gestion de Bases de Donnée.
 % \item[SI  :] \hspace{0.3cm}Système D’information. 
%  \item[SQL :] Standard Query language.
  \item[UML :] Unified Modeling Language.
  \item[HTML :] HyperText Markup Language
  \item[HTTP :] Hypertext Transfer Protocol
  \item[HTTPS :] Hypertext Transfer Protocol Secure
  \item[API :] Application Programming Interface
  \item[REST :] Representational State Transfer
  \item[URL :]  Uniform Resource Locator
  \item[UP :] Unified Process
  \item[2TUP :] two track unified process

 
  
  
\end{description}
%--------------------------liste des glossaire
%\include{glossaire}

\pagenumbering{arabic}
\setcounter{page}{1}
%--------------------------Introductions générale
\addcontentsline{toc}{chapter}{Introduction générale}
%\part*{Introduction}
\chapter*{Introduction générale}
\markboth{Introduction générale}D
Dans le domaine du BTP, ici en Algérie, il existe quelques moyens tels que les réseaux sociaux et des sites web tels que Ouedkniss, qui permet aux entreprises et aux individus de proposer leurs services et marchandises d’une manière anarchique et qui n’assure surtout pas l'évaluation des services et marchandises.
\\

\par Soit pour les grands projets qui nécessitent des entreprises pour les réaliser, soit pour les petits projets qui nécessitent un seul travailleur ou plusieurs, le client rencontre toujours des problème lors de la recherche de la bonne personne à recruter, car il existe pas une platform ou un organisme qui permet l’évaluation et recommandation de tell services. Un autre problème qu’on rencontre souvent aussi ici en Algérie, c’est le problème de prix, si vous cherchez les prix de n'importe quelle matière essentielle de batiment en ligne (en utilisant google par exemple), vous ne trouverez pas grand chose. Pour une personne qui ne connaît rien du BTP et des prix des matières du BTP, il peut se tromper facilement.\\

\par Notre projet est de concevoir une application web qui permet de résoudre cette problématique d’un côté et d'offrir des services de qualité soit pour les travailleurs ou les clients et permettre aussi d'éliminer l'ambiguïté sur le marché du BTP  ici en Algérie d’un autre côté.\\
\newpage
\vspace{0.5cm}
\par Notre mémoire se présente en un document structuré en quatre chapitres comme suit : 

\begin{itemize}[label=\textbullet]
\item \textbf{Chapitre 1 " Définition et généralité sur le système d'information" : } ------ \\

\item \textbf{Chapitre 2 " ******* " : } ------\\

\item \textbf{Chapitre 3 " Conception " :} ---------\\  

\item \textbf{Chapitre 4 " Réalisation" : }-------
\end{itemize}

\par Enfin, nous clôturons ce mémoire par une conclusion dans laquelle nous résumons notre mémoire et nous exposons quelques perspectives futures.























%--------------------------Différents chapitres
\include{chapitre1}
\documentclass[12pt, letterpaper]{article}
\usepackage[utf8]{inputenc}
\usepackage[T1]{fontenc}
\usepackage[french]{babel}



\begin{document}
\begin{center}
    \huge{Chapitre 2: Généralité sur le développement web}
\end{center}
\section{Introduction}


\section{}


\end{document}
\documentclass[12pt, letterpaper]{article}
\usepackage[utf8]{inputenc}
\usepackage[T1]{fontenc}
\usepackage[french]{babel}



\begin{document}
\begin{center}
    \huge{Chapitre 3: Conception}
\end{center}
\section{Introduction}


\section{}


\end{document}
\chapter{Réalisation }
\section{Introduction}


%--------------------------Conclusion générale
\addcontentsline{toc}{chapter}{Conclusion générale}
\chapter*{Conclusion générale}
\markboth{Conclusion générale}D

 
 
 
 
 


%--------------------------Références bibliographiques
\bibliography{refs}       
\bibliographystyle{plain} 

\end{document}